\chapter{Quick start}\label{ch:quickstart}

This chapter is a guided first session. In five steps we build a curve, turn it into a surface,
create a solid with a bore by a Boolean operation, compute its mass properties and write it to
a STEP file. A second script meshes the result with gmsh. Each step is explained briefly. The
following chapters cover the functions in depth.

The complete script is:

\example{03_quickstart.py}

Running it prints

\begin{lstlisting}[style=text]
arc: degree 2, 3 control points, rational True, domain (0.0, 1.0)
arc length 15.707963267949 (exact 15.707963267949)
arc(0.5) = (7.071068, 0.000000, 7.071068), |arc(0.5)| = 10.000000000000000
dome: degrees (2, 2), net 9 x 3
dome area 628.3185307180 (exact 628.3185307180)
part: 7 faces, closed True, labels ['back', 'bottom', 'front', 'left', 'right', 'side', 'top']
volume 14869.02664289 (exact 14869.02664471)
centroid (0.000000, -0.000000, 5.000000)
inside at (0, 0, 5): False, at (10, 10, 5): True
read back: 1 shell(s), volume 14869.02664989
\end{lstlisting}

\section{Step 1: a curve}\label{sec:qs-curve}

\begin{lstlisting}
arc = nf.make_arc(center=(0, 0, 0), x_axis=(1, 0, 0), y_axis=(0, 0, 1),
                  radius=10.0, start_angle=0.0, end_angle=math.pi / 2)
\end{lstlisting}

\apiref{make_arc} creates a circular arc. The angles are measured in radians from
\code{x_axis} towards \code{y_axis}, so this arc runs from $(10,0,0)$ to $(0,0,10)$ in the
$xz$ plane. Points are given as tuples. Any sequence of three numbers works, and so does an
\apiref{Vec3}.

The result is an \apiref{NurbsCurve} of degree 2 with three control points and weights
$1, \tfrac12\sqrt2, 1$: the classical exact representation of a circular arc as a rational
quadratic. This exactness shows in the output:
\begin{itemize}
  \item \apiref{length} integrates the arc length adaptively and agrees with $5\pi$ to all
    printed digits;
  \item \code{arc(0.5)} evaluates the curve. Calling a curve is the same as
    \apiref{NurbsCurve.evaluate}, and the point has distance exactly 10 from the centre.
\end{itemize}

Every curve created by a constructor has the parameter domain $[0,1]$ (\code{arc.domain}). The
parameter is \emph{not} proportional to arc length; for a rational arc the midpoint happens to
coincide only because of symmetry. \Cref{sec:concepts-nurbs} explains why.

\section{Step 2: a surface}\label{sec:qs-surface}

\begin{lstlisting}
dome = nf.revolve(arc, axis_point=(0, 0, 0), axis_dir=(0, 0, 1), angle=2 * math.pi)
\end{lstlisting}

\apiref{revolve} sweeps the profile around an axis. A full turn of a quarter circle about the
$z$ axis is a hemisphere of radius 10. The result is an exact \apiref{NurbsSurface}: the
circumferential direction $u$ is a rational quadratic circle with 9 control points, and the
profile direction $v$ keeps the 3 control points of the arc. Its area, computed by
\apiref{area}, agrees with $2\pi R^2 = 200\pi$. Surfaces are evaluated with \code{dome(u, v)}.
\Cref{ch:surfaces} describes all ways to build surfaces.

\section{Step 3: solids and a Boolean operation}\label{sec:qs-solid}

\begin{lstlisting}
box = nf.make_box(corner=(-20, -20, 0), size=(40, 40, 10))
bore = nf.make_cylinder(base=(0, 0, -5), axis=(0, 0, 1), radius=6.0, height=20.0)
part = nf.boolean(box, bore, nf.BooleanOp.Difference)
\end{lstlisting}

Solids in \NF{} are closed \apiref{Shell} objects: sets of trimmed surfaces (\emph{faces}) sewn
together along shared edges, with normals pointing outwards. \apiref{make_box} and
\apiref{make_cylinder} create such shells directly. The cylinder is longer than the box is high,
so it cuts through cleanly. Coincident faces are a degenerate case for any Boolean algorithm and
should be avoided.

\apiref{boolean} with \code{BooleanOp.Difference} subtracts the cylinder from the box. The result
has 7 faces: the six box faces (top and bottom now with a circular hole) and the cylindrical
bore wall. Every face carries a \emph{label}. The primitives name their faces (\code{bottom},
\code{top}, \code{front}, \code{back}, \code{left}, \code{right} for the box; \code{side},
\code{bottom}, \code{top} for the cylinder), and Booleans keep the labels. The bore wall is
therefore \code{side}. Labels later become physical groups in the mesh, where boundary conditions
are assigned (\cref{ch:solids,ch:meshing}).

\section{Step 4: mass properties and point classification}\label{sec:qs-mass}

\begin{lstlisting}
mp = part.mass_properties()
nf.point_in_shell(part, (0, 0, 5))
\end{lstlisting}

\apiref{Shell.mass_properties} integrates over the faces with the divergence theorem and returns
volume, centroid, the inertia tensor about the centroid (unit density) and the principal moments
and axes. The volume agrees with the exact value $40\cdot40\cdot10 - 36\pi\cdot 10$ to a relative
$1.2\cdot10^{-10}$. That is better than the default integration tolerance \code{rel_tol=1e-8},
which bounds the error. \apiref{point_in_shell} classifies points by ray casting: the axis of
the bore lies outside the material.

\section{Step 5: STEP export and import}\label{sec:qs-step}

\begin{lstlisting}
nf.write_step("part.step", part, unit="mm", name="drilled block")
back = nf.read_step("part.step")
\end{lstlisting}

\apiref{write_step} uses \NF{}'s own STEP writer. No CAD library is involved. The file declares
millimetres as its length unit and contains one solid with exact B-spline faces. Any CAD system
can open it. \apiref{read_step} reads STEP files from other systems or from \NF{} and returns a
list of shells. The volume of the re-imported part differs from the original by a relative
$5\cdot10^{-10}$. The reader rebuilds the trim loops from the 3D edges, so the faces of the
read-back shell are trimmed by slightly different (but equivalent) curves. \Cref{ch:step} covers units, names and the supported entities.

\section{Step 6: meshing}\label{sec:qs-mesh}

Meshing needs gmsh (\code{pip install "nurbforge[mesh]"}). The first line of the script tells
the test suite to skip it when gmsh is not available:

\example{03_quickstart_mesh.py}

\begin{lstlisting}[style=text]
face labels -> gmsh surfaces: {'bottom': [1], 'top': [2], 'front': [3], 'back': [4], 'left': [5], 'right': [6], 'side': [7]}
BEM mesh: 796 nodes, 1592 triangles, watertight True, volume 14896.091 (exact 14869.027)
\end{lstlisting}

\begin{itemize}
  \item \apiref{mesh.GmshSession} wraps the gmsh session. \apiref{mesh.GmshSession.add_shell}
    hands the solid over with its exact NURBS faces and, with \code{physical_groups=True}, one
    named physical surface group per face label plus a physical volume.
  \item \apiref{mesh.GmshSession.refine} sets a smaller element size on the faces labelled
    \code{side} (the bore), and \apiref{mesh.GmshSession.mesh} generates a tetrahedral mesh
    (\code{dim=3}) with a global size of 4.
  \item \apiref{mesh.bem_mesh} produces a closed triangle surface mesh whose normals point out of
    the solid, as boundary element solvers need it, and returns it as NumPy arrays. Its
    volume is slightly larger than the exact one: the polygonal bore removes less material
    than the true cylinder. The difference shrinks with the element size.
\end{itemize}

\section{Where to go next}\label{sec:qs-next}

\begin{itemize}
  \item \Cref{ch:concepts} explains NURBS, parameters, tolerances and errors. Read it before
    you rely on numbers.
  \item \Cref{ch:curves,ch:surfaces} cover all curve and surface constructions,
    \cref{ch:solids} shells, solids and Booleans.
  \item \Cref{ch:meshing} shows local mesh sizes, high-order elements and BEM meshes in
    detail, \cref{ch:step} the STEP interface.
  \item The repository contains further scripts in \code{examples/python/}, from lens design to
    slicing a microlens for two-photon lithography.
\end{itemize}
